\documentclass[10pt,a4paper]{amsart}
\usepackage{amsmath,amssymb,amsthm,graphicx}
\usepackage{verbatim}
\newenvironment{reference}{\comment}{\endcomment}
%\newenvironment{reference}{}{}
\newenvironment{slogan}{\comment}{\endcomment}
\newenvironment{history}{\comment}{\endcomment}

% For commutative diagrams we use Xy-pic
\usepackage[all]{xy}

% We use 2cell for 2-commutative diagrams.
\xyoption{2cell}
% Theorem environments.
%
\theoremstyle{plain}
\newtheorem{theorem}[subsection]{Theorem}
\newtheorem{proposition}[subsection]{Proposition}
\newtheorem{lemma}[subsection]{Lemma}

\theoremstyle{definition}
\newtheorem{definition}[subsection]{Definition}
\newtheorem{example}[subsection]{Example}
\newtheorem{exercise}[subsection]{Exercise}
\newtheorem{situation}[subsection]{Situation}

\theoremstyle{remark}
\newtheorem{remark}[subsection]{Remark}
\newtheorem{remarks}[subsection]{Remarks}

\numberwithin{equation}{subsection}

% Macros
%
\def\lim{\mathop{\mathrm{lim}}\nolimits}
\def\colim{\mathop{\mathrm{colim}}\nolimits}
\def\Spec{\mathop{\mathrm{Spec}}}
\def\Hom{\mathop{\mathrm{Hom}}\nolimits}
\def\Ext{\mathop{\mathrm{Ext}}\nolimits}
\def\SheafHom{\mathop{\mathcal{H}\!\mathit{om}}\nolimits}
\def\SheafExt{\mathop{\mathcal{E}\!\mathit{xt}}\nolimits}
\def\Sch{\mathit{Sch}}
\def\Mor{\mathop{\mathrm{Mor}}\nolimits}
\def\Ob{\mathop{\mathrm{Ob}}\nolimits}
\def\Sh{\mathop{\mathit{Sh}}\nolimits}
\def\NL{\mathop{N\!L}\nolimits}
\def\CH{\mathop{\mathrm{CH}}\nolimits}
\def\proetale{{pro\text{-}\acute{e}tale}}
\def\etale{{\acute{e}tale}}
\def\QCoh{\mathit{QCoh}}
\def\Ker{\mathop{\mathrm{Ker}}}
\def\Im{\mathop{\mathrm{Im}}}
\def\Coker{\mathop{\mathrm{Coker}}}
\def\Coim{\mathop{\mathrm{Coim}}}

% Boxtimes
%
\DeclareMathSymbol{\boxtimes}{\mathbin}{AMSa}{"02}

%
% Macros for moduli stacks/spaces
%
\def\QCohstack{\mathcal{QC}\!\mathit{oh}}
\def\Cohstack{\mathcal{C}\!\mathit{oh}}
\def\Spacesstack{\mathcal{S}\!\mathit{paces}}
\def\Quotfunctor{\mathrm{Quot}}
\def\Hilbfunctor{\mathrm{Hilb}}
\def\Curvesstack{\mathcal{C}\!\mathit{urves}}
\def\Polarizedstack{\mathcal{P}\!\mathit{olarized}}
\def\Complexesstack{\mathcal{C}\!\mathit{omplexes}}
% \Pic is the operator that assigns to X its picard group, usage \Pic(X)
% \Picardstack_{X/B} denotes the Picard stack of X over B
% \Picardfunctor_{X/B} denotes the Picard functor of X over B
\def\Pic{\mathop{\mathrm{Pic}}\nolimits}
\def\Picardstack{\mathcal{P}\!\mathit{ic}}
\def\Picardfunctor{\mathrm{Pic}}
\def\Deformationcategory{\mathcal{D}\!\mathit{ef}}
\usepackage{cleveref}
\crefname{originallemma}{original lemma}{original lemmas}
\Crefname{originallemma}{Original lemma}{Original lemmas}
\makeatletter
\newlabel{lemma-equivalence-coherent}{{0GPK}{source}{}{sourcepointer-1}{}}
\newlabel{lemma-equivalence-coherent@cref}{{[originallemma][0][]0GPK}{[1][0][]source}{}{}{}}
\makeatother
\crefname{originalremark}{original remark}{original remarks}
\Crefname{originalremark}{Original remark}{Original remarks}
\makeatletter
\newlabel{remark-equivalence-coherent-linear}{{0GPL}{source}{}{sourcepointer-2}{}}
\newlabel{remark-equivalence-coherent-linear@cref}{{[originalremark][0][]0GPL}{[1][0][]source}{}{}{}}
\makeatother
\crefname{originalreference}{original reference}{original references}
\Crefname{originalreference}{Original reference}{Original references}
\makeatletter
\newlabel{Gabriel}{{shared-bibliography-key}{source}{}{sourcepointer-3}{}}
\newlabel{Gabriel@cref}{{[originalreference][0][]shared-bibliography-key}{[1][0][]source}{}{}{}}
\makeatother
\title{A k-linear equivalence of coherent categories over finite type schemes comes from a k-isomorphism and an invertible module (weak form of Gabriel)}
\author{The Stacks Project Authors}
\date{stacks-project:functors.tex:local-snapshot-20261012; GNU FDL 1.2 or later}
\begin{document}
\maketitle
\noindent\textit{Editorial note.} All delivered mathematical text is the author's own, copied verbatim from the pinned Stacks Project chapter file functors.tex; no statement or proof step was rewritten, reordered or removed.\par\medskip
\noindent\textit{Editorial note.} Span recomputation by label pairing (LANE\_BRIEF 10.2): in this chapter the author's proof environment follows the statement environment as a standalone block; the spec covers the complete environment pair (statement span, proof span) as recorded in the pool candidate. This item is a dual-proof environment: the spec-level proof\_lines override collects both proof environments (LANE\_BRIEF 10.2 span recomputation); the pool candidate recorded only the first.\par\medskip
\noindent\textit{Editorial note.} Occupation scan: combined claim\_key + statement\_anchor + generic-line scan over all registration files (active and withdrawn) found 0 hits for functors.tex (fu-occupied.json, 2026-10-14); consistent with the duty manager's six-check dispatch.\par\medskip
\ifdefined\hypertarget\hypertarget{sourcepointer-1}{}\fi
\noindent\textit{Source reference.} Original lemma 0GPK: sources/stacks-chow/tags.txt, source lines 20830--20830. Cross-chapter standard prerequisite cited by the author inside the counted span; it is used as a pointer and its own proof is not part of this item.\par\medskip
\ifdefined\hypertarget\hypertarget{sourcepointer-2}{}\fi
\noindent\textit{Source reference.} Original remark 0GPL: sources/stacks-chow/tags.txt, source lines 20831--20831. Cross-chapter standard prerequisite cited by the author inside the counted span; it is used as a pointer and its own proof is not part of this item.\par\medskip
\ifdefined\hypertarget\hypertarget{sourcepointer-3}{}\fi
\noindent\textit{Source reference.} Original reference shared-bibliography-key: sources/stacks-functors/SOURCE.json, source lines 14--15. Author's bibliographic citation inside the counted span (reference environment pointing at Gabriel); the shared Stacks bibliography is not staged locally, so the cite key is declared here per SOURCE.json bibliography\_note. Pointer only; its content is not part of this item.\par\medskip
\begin{lemma}
\label{lemma-equivalence-coherent-over-field}
\begin{reference}
Weak version of the result in \cite{Gabriel}
stating that the category of quasi-coherent modules
determines the isomorphism class of a scheme.
\end{reference}
Let $k$ be a field. Let $X$, $Y$ be finite type schemes over $k$ with
$X$ separated and $Y$ reduced. If there is a $k$-linear equivalence
$F : \textit{Coh}(\mathcal{O}_X) \to \textit{Coh}(\mathcal{O}_Y)$
of categories, then there is an isomorphism $f : Y \to X$
over $k$ and an invertible $\mathcal{O}_Y$-module $\mathcal{L}$
such that $F(\mathcal{F}) = f^*\mathcal{F} \otimes \mathcal{L}$.
\end{lemma}

\begin{proof}[Proof using Gabriel-Rosenberg reconstruction]
This lemma is a weak form of the results discussed in
Lemma \ref{lemma-equivalence-coherent} and
Remark \ref{remark-equivalence-coherent-linear}.
\end{proof}

\end{document}
